\documentclass[letter, 12pt]{article}
\usepackage[margin=1in, paperwidth=8.5in, paperheight=11in]{geometry}
\usepackage{ifpdf,amsmath, amssymb, comment, color, graphicx, stmaryrd,setspace,enumerate,tikz, fancyhdr, wrapfig, textcomp, units, mathptmx, multicol}
\hyphenpenalty=5000

\newcommand{\Q}{\mathbb{Q}}
\newcommand{\R}{\mathbb{R}}
\newcommand{\Z}{\mathbb{Z}}

\newenvironment{red}{ %\iffalse
\color{red}}{\ignorespacesafterend
} %\fi
\newcommand{\blue}[1]{\textcolor{blue}{#1}}
\newcommand{\green}[1]{\textcolor{green}{#1}}
\renewcommand{\section}[1]{\begin{center} \textbf{#1} \\\end{center}}
\newcommand{\diff}[1]{\frac{d}{d#1}}
%
\setlength{\parindent}{0in}
%\oddsidemargin=-.25in
\allowdisplaybreaks
\pagestyle{fancy}
\renewcommand{\headrulewidth}{0pt}
\lhead{Calculus II}
\rhead{Spring 2024}
%\lfoot{\copyright\ CLEAR Calculus 2010}
\cfoot{}
\renewcommand{\thefootnote}{*}

\newcommand{\itx}[1]{\intertext{\begin{flushright}#1\end{flushright}}}

\begin{document}
\section{Derivatives Review Problems -- \begin{red}
Solutions!
\end{red}}

\textbf{Instructions:} 
In this Problem Set, you'll carefully solve one of these differentiation problems, then make a purposeful mistake somewhere. Follow these steps:
\begin{enumerate}[(a)]
	\item Choose one of the problems below.
	\item Refer to the list on the next page to identify \textbf{all} the derivative rules you're going to need to use for your problem.
	\item Carefully find the derivative \textbf{step by step}, write out \textbf{every step}, and list all the rules you used in each step. Please see the example on the next page to see how detailed your steps should be.
	\item Choose a \textbf{second} problem, and do the same things, but make your answer wrong somehow. Bonus points\footnote{Points don't exist.} for an error you think people probably make in real life.
\end{enumerate}
Throughout, carefully label any derivative you find by name, using correct derivative notation.

\hrulefill

\begin{enumerate}[{\bf (1)}]
\item
Find $\frac{dy}{dx}$ if \(y=(2x-1)^8-x+e\).

\begin{red}
	\begin{multicols}{2}
		\begin{itemize}
			\item constant multiple rule
			\item sum and difference rule
			\item power rule
			\item chain rule
		\end{itemize}
	\end{multicols}
	
	\begin{align*}
	\frac{dy}{dx} &= \diff{x}(2x-1)^8 - \diff{x}(x) + \diff{x}(e)
	& \textrm{(sum and difference rule)} \\
	&= 8(2x-1)^7\cdot\diff{x}(2x-1) - 1 + 0 
	& \textrm{(chain rule, power rule)} \\
	&= 8(2x-1)^7 \cdot (2 - 0) - 1 
	& \textrm{(constant multiple rule, power rule)}
	\end{align*}
\end{red}
 
\item 
If \(L(y) = (\cos(y))^{3} + \cos(3)\), find $L'(y)$.

\begin{red}
	\begin{multicols}{2}
		\begin{itemize}
			\item sum and difference rule
			\item trig functions rule
			\item power rule
			\item chain rule
		\end{itemize}
	\end{multicols}
	
	\begin{align*}
	L'(y) &= \diff{y}(\cos(y))^3 + \diff{y}\cos(3)
	& \textrm{(sum and difference rule)} \\
	&= 3(\cos(y))^2\cdot\diff{y}\cos(y) + 0
	& \textrm{(chain rule, power rule)} \\
	&= 3(\cos(y))^2\cdot(-\sin(y))
	& \textrm{(trig functions rule)}
	\end{align*}
\end{red}
 
\pagebreak

\item
Find $B'(x)$ if \(B(x) = \cos(e^{3 x})\).

\begin{red}
	\begin{multicols}{2}
		\begin{itemize}
			\item chain rule
			\item trig functions rule
			\item exponential functions rule
			\item constant multiple rule
			\item power rule
			\item[]
		\end{itemize}
	\end{multicols}
	
	\begin{align*}
	B'(x) &= -\sin(e^{3x})\cdot\diff{x}(e^{3x})
	& \textrm{(chain rule, trig functions rule)} \\
	&= -\sin(e^{3x})\cdot e^{3x} \cdot \diff{x}(3x) 
	& \textrm{(chain rule, exponential functions rule)} \\
	&= -\sin(e^{3x})\cdot e^{3x} \cdot 3
	& \textrm{(constant multiple rule, power rule)}
	\end{align*}
\end{red} 
 
\item
Differentiate \(W(x) =\displaystyle{x^4 - 6x + 10 \over \sqrt{7x + 8}}\).

\begin{red}
	\begin{multicols}{2}
		\begin{itemize}
			\item quotient rule
			\item sum and difference rule
			\item power rule
			\item constant multiple rule
			\item chain rule
			\item[]
		\end{itemize}
	\end{multicols}
	
	\begin{align*}
	W'(x) &= \frac{(\sqrt{7x+8})\cdot\diff{x}(x^4-6x+10) 
	               - (x^4-6x+10)\cdot\diff{x}(\sqrt{7x+8})}
	              {(\sqrt{7x+8})^2} \\
	\itx{(quotient rule)} 
	&= \frac{(\sqrt{7x+8})\cdot(4x^3-6) 
	               - (x^4-6x+10)\cdot(\tfrac{1}{2}(7x+8)^{-1/2}\cdot 7) }
	              {(\sqrt{7x+8})^2} \\
	\itx{(power rule, sum and difference rule, constant multiple rule, chain rule)}
	\end{align*}
\end{red}
\pagebreak

\item Differentiate \(g(q) = \displaystyle{q^{5} \over \ln(q + {1\over3})}\).
\begin{red}
	\begin{multicols}{2}
		\begin{itemize}
			\item quotient rule
			\item power rule
			\item chain rule
			\item logarithmic functions rule
			\item sum and difference rule
			\item[]
		\end{itemize}
	\end{multicols}
	\begin{align*}
	g'(q) &= \frac{(\ln(q+\tfrac{1}{3}))\cdot\diff{q}(q^5) - 
	              (q^5)\cdot\diff{q}(\ln(q+\tfrac{1}{3}))}
	              {(\ln(q+\tfrac{1}{3}))^2}
	\itx{(quotient rule)} 
	&= \frac{(\ln(q+\tfrac{1}{3}))\cdot(5q^4) - 
	              (q^5)\cdot\left(\frac{1}{q+\tfrac{1}{3}}\cdot (1+0)\right)}
	              {(\ln(q+\tfrac{1}{3}))^2} \\
	\itx{(power rule, chain rule, log. functions rule, sum and difference rule)}
	\end{align*}
\end{red}


\item
Differentiate \(B(x) = (C e^{3x} - D)(2 x)\), where \(C\) and \(D\) are constants.

\begin{red}
    	\begin{multicols}{2}
		\begin{itemize}
			\item product rule
			\item sum and difference rule
			\item constant multiple rule
			\item exponential functions rule
			\item power rule
			\item[]
		\end{itemize}
	\end{multicols}
	
	\begin{align*}
	B'(x) &= \left(\diff{x}(C e^{3x} - D)\right)\cdot(2 x) + (C e^{3x} - D)\cdot\diff{x}(2 x) \\
	\itx{(product rule)}
	&= \left(C\cdot\diff{x}(e^{3x}) - \diff{x}(D)\right) \cdot(2x) + (C e^{3x} - D)\cdot(2) \\
	\itx{(sum and difference rule, constant multiple rule, power rule)} \\
	&= \left(C\cdot(e^{3x}\cdot(3)) - 0\right) \cdot(2x) + (C e^{3x} - D)\cdot(2) \\
	\itx{(chain rule, exponential functions rule, constant multiple rule, power rule)}
	\end{align*}
\end{red}

\item
Differentiate \(f(y) = \left(\sqrt{y^4 + 1}\right) \left(\ln\left(\frac{y}{2}\right)\right)\).

\begin{red}
	\begin{multicols}{2}
		\begin{itemize}
			\item product rule
			\item chain rule
			\item logarithmic functions rule
			\item constant multiple rule
			\item power rule
			\item sum and difference rule
			\item \textbf{NOT the quotient rule!!}
			\item[]
		\end{itemize}
	\end{multicols}
	\begin{align*}
	f'(y) &= \left(\sqrt{y^4 + 1}\right)\cdot \diff{y} \left(\ln\left(\frac{y}{2}\right)\right)
	      + \left[\diff{y}\left(\sqrt{y^4 + 1}\right)\right]\cdot 
	      			\left(\ln\left(\frac{y}{2}\right)\right) \\
	\itx{(product rule)}
	&= \left(\sqrt{y^4 + 1}\right)\cdot \left(\frac{1}{y/2}\cdot\diff{y}\left(\frac{y}{2}\right)\right)
	      + \left[\frac{1}{2}(y^4 + 1)^{-1/2}\cdot\diff{y}(y^4+1)\right]\cdot 
	      			\left(\ln\left(\frac{y}{2}\right)\right) \\
	\itx{(chain rule, logarithmic function rule, power rule)}
	&= \left(\sqrt{y^4 + 1}\right)\cdot \left(\frac{1}{y/2}\cdot\left(\frac{1}{2}\right)\right)
	      + \left[\frac{1}{2}(y^4 + 1)^{-1/2}\cdot(4y^3+0)\right]\cdot 
	      			\left(\ln\left(\frac{y}{2}\right)\right) \\
	\itx{(constant multiple rule, power rule, sum and difference rule)}
	\end{align*}

\end{red}
\pagebreak

\item
Find $\frac{dB}{dx}$ if 
\(B =\displaystyle{e^{x + \pi} \over \sin(\pi x)}\).
\begin{red}
    \begin{multicols}{2}
		\begin{itemize}
			\item quotient rule
			\item chain rule
			\item exponential functions rule
			\item sum and difference rule
			\item power rule
			\item trig functions rule
			\item constant multiple rule
			\item[]
		\end{itemize}
	\end{multicols}
	\begin{align*}
	\frac{dB}{dx} &= \frac{(\sin(\pi x))\cdot\diff{x}(e^{x + \pi}) 
	               - (e^{x + \pi})\cdot\diff{x}(\sin(\pi x))}
	              {(\sin(\pi x))^2} \\
	\itx{(quotient rule)}
	&= \frac{(\sin(\pi x))\cdot(e^{x + \pi}\cdot(1+0)) 
	               - (e^{x + \pi})\cdot(\cos(\pi x)\cdot\pi)}
	              {(\sin(\pi x))^2} \\
	\itx{(chain rule, exponential functions rule, sum and difference rule, \\ 
	power rule, trig functions rule, constant multiple rule)}
	\end{align*}
\end{red}



\item 
Differentiate \(h(p) = \displaystyle{(2 p - 1)^{1/4} \over \cos(p + 5)}\).

\begin{red}
	\begin{multicols}{2}
		\begin{itemize}
			\item quotient rule
			\item chain rule
			\item power rule
			\item constant multiple rule
			\item sum and difference rule
			\item trig functions rule
		\end{itemize}
	\end{multicols}
	\begin{align*}
	h'(p) &= \frac{\cos(p + 5)\cdot\diff{p}((2 p - 1)^{1/4}) 
	               - ((2 p - 1)^{1/4})\cdot\diff{p}(\cos(p + 5))}
	              {(\cos(p + 5))^2} \\
	\itx{(quotient rule)}
	&= \frac{\cos(p + 5)\cdot\left[\frac{1}{4}(2 p - 1)^{-3/4}\cdot(2-0)\right] 
	               - ((2 p - 1)^{1/4})\cdot\diff{p}(\cos(p + 5))}
	              {(\cos(p + 5))^2} \\
	\itx{(chain rule, power rule, sum and difference rule)}
	&= \frac{\cos(p + 5)\cdot\left[\frac{1}{4}(2 p - 1)^{-3/4}\cdot(2-0)\right] 
	               - ((2 p - 1)^{1/4})\cdot(-\sin(p + 5)\cdot 1)}
	              {(\cos(p + 5))^2} \\
	\itx{(chain rule, trig functions rule, power rule)}
	\end{align*}
\end{red}

\item
Differentiate \(f(q)=(\cos q)\cdot (\sin (\tfrac{2}{\pi} q)) \).

\begin{red}
    \begin{multicols}{2}
		\begin{itemize}
			\item product rule
			\item trig functions rule
			\item chain rule
			\item constant multiple rule
			\item power rule
			\item
		\end{itemize}
	\end{multicols}
	\begin{align*}
	f'(q) &= \left[\diff{q}(\cos(q))\right] \cdot (\sin (\tfrac{2}{\pi} q))
		+ (\cos(q)) \cdot \diff{q} (\sin (\tfrac{2}{\pi} q)) \\
	\itx{(product rule)}
	&= (-\sin(q)) \cdot (\sin (\tfrac{2}{\pi} q))
		+ (\cos(q)) \cdot (\cos (\tfrac{2}{\pi} q) \cdot \tfrac{2}{\pi}) \\
	\itx{(trig functions rule, chain rule, constant multiple rule, power rule)}		
	\end{align*}
\end{red}

\pagebreak

\item
Differentiate \(G = p x^5 + q x^3 - x^r\), where \(p,q\) and \(r\)
are constants.
\begin{red}
    	\begin{multicols}{2}
		\begin{itemize}
			\item sum and difference rule
			\item constant multiple rule
			\item power rule
		\end{itemize}
	\end{multicols}
	\begin{align*}
	\frac{dG}{dx} &= p\cdot\diff{x}(x^5) + q\cdot\diff{x}(x^3) - \diff{x}(x^r)
	& \textrm{(sum and difference rule, constant multiple rule)}\\
	&&\\
	&= 5p x^4 + 3q x^2 - rx^{r-1} 
	& \textrm{(power rule)}
	\end{align*}
\end{red}

\item
Differentiate \(V = (3 x - \tfrac{1}{4})^{-4} - (5x)^{5/4}\).
\begin{red}
	\begin{multicols}{2}
		\begin{itemize}
			\item sum and difference rule
			\item chain rule
			\item power rule
			\item constant multiple rule
		\end{itemize}
	\end{multicols}
	\begin{align*}
	\frac{dV}{dx} &= \diff{x}\left[(3 x - \tfrac{1}{4})^{-4}\right] - \diff{x}\left[(5x)^{5/4}\right] \\
	\itx{(sum and difference rule)}
	&= -4(3 x - \tfrac{1}{4})^{-5} \cdot (3-0) - \tfrac{5}{4} (5x)^{1/4} \cdot 5 \\
	\itx{(chain rule, power rule, sum and difference rule, constant multiple rule)}
	\end{align*}
\end{red}
\pagebreak

\item
Differentiate \(h(x) = \sqrt{e^{6 x}}\).
\begin{red}
    \begin{multicols}{2}
		\begin{itemize}
			\item chain rule
			\item power rule
			\item exponential functions rule
			\item constant multiple rule
		\end{itemize}
	\end{multicols}
	\begin{align*}
	h'(x) &= \tfrac{1}{2}(e^{6x})^{-1/2} \cdot\diff{x}(e^{6x}) &
	\textrm{(chain rule, power rule)}\\
	&= \tfrac{1}{2}(e^{6x})^{-1/2} \cdot (e^{6x} \cdot \diff{x}(6x)) &
	\textrm{(chain rule, exponential functions rule)}\\
	&= \tfrac{1}{2}(e^{6x})^{-1/2} \cdot (e^{6x} \cdot 6) &
	\textrm{(constant multiple rule, power rule)}\\
	\end{align*}
	(\textbf{Note:} There's another way to do this one. Note that $h(x) = \left[e^{6x}\right]^{1/2}
	= e^{\frac{6}{2}x} = e^{3x}$. \\
	 Then $h'(x) = 3e^{3x}$. Check that this is the same as the other answer above!
\end{red}
	
\item
Differentiate \(h(t) = -3(2t^4 + 7)^4 + 3\).
\begin{red}
    \begin{multicols}{2}
		\begin{itemize}
			\item sum and difference rule
			\item constant multiple rule
			\item chain rule
			\item power rule
		\end{itemize}
	\end{multicols}
	\begin{align*}
	h'(t) &= -3\diff{t}(2t^4 + 7)^4 + \diff{t}(3) \\
	\itx{(sum and difference rule, constant multiple rule)}
	&= -3\left[4(2t^4+7)^3\cdot\diff{t}(2t^4+7)\right] + 0 \\
	\itx{(chain rule, power rule)}
	&= -3\left[4(2t^4+7)^3\cdot(8t^3 + 0)\right] \\
	\itx{(sum and difference rule, constant multiple rule, power rule)}
	\end{align*}
\end{red}
	
\end{enumerate}


\end{document}
